\subsection{Predicting the new games}
\label{sec:predict_new_games}

We now introduce the four novel strategic games.
To the best of our knowledge, these games are not in the LLM's training corpus, thus providing a stringent testing ground for the \aisub{s} we have explored so far.
Three of the games parallel AR's games in strategic structure but modify the implementation: participants choose between 1 and 10 (rather than 11 and 20) and earn points instead of shekels.
The instructions for the ``basic'' version of this 1-10 game highlight these differences:
\begin{quote}
\vspace{-.5cm}
\singlespacing
\emph{You are going to play a game where you must select a whole number between 1 and 10.
You will receive a number of points equivalent to that number.
After you tell us your number, we will randomly pair you with another player who is also playing this game. They will also have chosen a number between 1 and 10.
If either of you selects a number exactly one less than the other player's number, the player with the lower number will receive an additional 10 points.}
\end{quote}
We adapted the costless and cycle variants similarly.
The fourth ``1-7 game'' introduces an entirely new variant with a restricted choice set (see Appendix~\ref{app:instruct} for the full instructions for all games). 

Of the 1,000 participants we recruited from Prolific, \NTotal{} passed the validation check and were randomly distributed across the four games.
To ensure incentive compatibility, participants were paid \$1 for completion and had a 10\% chance of receiving the dollar value of the points they earned.
We preregistered our complete experimental design, including all prompts for both the baseline and selected \aisub{s}---the latter using only the weights optimized on the basic 11-20 game.
We have $\bm{\theta}^*$, $\bm{\theta}^*_{Hist}$, $\bm{\theta}^*_{MB}$, $\bm{\theta}^*_{N}$, and the baseline play each game.
All \aisub{} responses are elicited using \textsc{Gpt-4o} with the temperature set to 1.
All \aisub{} samples played these games \emph{before} the human subjects' data was collected.

Figure~\ref{fig:all_110_a} presents the response distributions of initial play for all subject samples---both human (top two rows) and AI (remaining rows)---across the four novel games. 
The responses from AR are all shifted down by 10 for ease of comparison---20 is now 10, 19 is 9, etc.
Response distributions from the Prolific sample (thin black bars) are superimposed on all other samples.

\begin{figure}[h!]
\begin{center}
\caption{Analysis of novel 1-10 games: response distributions and KL divergences}
\vspace*{-0.1in}
\label{fig:all_110}
\begin{subfigure}{\textwidth}  
\refstepcounter{subfigure}\label{fig:all_110_a}
\centering
\begin{overpic}[width=\textwidth]{plots/all_110.pdf}
\put(-1,92){\large\bfseries (a)}
\end{overpic}
\end{subfigure}
\begin{subfigure}{\textwidth}
\refstepcounter{subfigure}\label{fig:all_110_b}
\centering
\begin{overpic}[width=\textwidth]{plots/kl_110.pdf}
\put(-1,23){\large\bfseries (b)}
\end{overpic}
\end{subfigure}
\end{center}
\vspace*{-0.2in}
\begin{footnotesize}
\begin{singlespace}
\textit{Notes:} (a) plots human and AI response distributions for the four games. 
Prolific data are in black superimposed on all other distributions. \citeauthor{1120Arad2012} data occupy the second row, shifted down by 10 to fit the 1-10 format.
The remaining rows show the various \aisub{} samples named in the right-hand column. 
(b) Reports the KL divergence between Prolific data and each other distribution.
The minimum in each panel is flagged by a checkmark.
\end{singlespace}
\end{footnotesize}
\end{figure}

The Prolific responses differ notably from AR's original results.
Specifically, in the basic 1-10 game, the Prolific sample's distribution is more uniform, with a modal choice at 8 rather than AR's mode at 7; in the costless variant, Prolific responses peak at 10 instead of AR's 9 and 8; and the cycle variant yields a more uniform distribution relative to AR's original sample. 
The newly introduced 1-7 game lacks an analogous comparison in AR, but the modal response is 5, and most other choices are on 4, 6, and 7.

Assuming the difference is not due to sampling variation, the gap between the Prolific data and \citeauthor{1120Arad2012} is somewhat surprising. 
Both sets of games share the same fundamental strategic structure. 
Although the games differ in their Nash equilibria (which we know is not how actual humans play anyway), a fixed set of $k$-level reasoners would play them identically: level-0 would naturally choose the highest number (10), level-1 the second-highest (9), and so on. 
The divergence may therefore reflect setting-specific factors or differences in beliefs about others' reasoning---precisely the kinds of factors that otherwise make theory difficult to apply.
In the end, even highly relevant human data proves to be an imperfect predictor.

Moving down Figure~\ref{fig:all_110_a}, the remaining rows show response distributions from various AI samples.
Whether or not the AI sample successfully passed validation on the costless and cycle games is indicated below the sample name.
The baseline AI (red) generally provides a poor fit to the Prolific data, frequently concentrating choices on 9. 
The notable exception is the 1-7 game, where it disperses responses across several choices.

Critically, the optimized sample of strategically-motivated \aisub{s} robustly generalizes to these novel settings.
This is the only sample of \aisub{s} that was successfully validated on all of the data from AR's original games (the atheoretical samples each failed to improve in at least one of these games).
Figure~\ref{fig:all_110_b} reports the KL divergence between the Prolific data and each AI-generated distribution: the strategic agents consistently outperform the baseline AI by at least \MinPercImproveArad\% in every variant (Appendix Figure~\ref{fig:dist_compare_1_10} shows the same comparisons for several alternative distance metrics). 
Especially notable is the strategic agents' near-perfect alignment in the entirely novel 1-7 game (KL divergence = \KlHumanOptDistSeven). 
Indeed, the strategically motivated AI even outperforms AR's original human data in predicting our Prolific participants' responses in the basic and cycle games.

In stark contrast, atheoretical \aisub{s} fail to generalize.
All arbitrary samples predict the human responses no better than the baseline in at least two of the four games.
The  ``Always Pick `N''' agents are particularly nonsensical as they are solely instructed to select integers between 11 and 20, 
highlighting a severe case of overfitting to a particular data-generating process.

Overall, these results underscore our central claim: carefully identifying theoretically-grounded candidate \persona{s} and validating their predictive utility in related but distinct contexts can substantially enhance predictive accuracy in novel, unseen settings.
Only agents subject to this approach generalized to the novel strategic games. 
